\section{코드 분석과 공동 개발 가이드}

본 장은 앞선 분석(1\textasciitilde 12장)을 실제 \emph{공동 개발}로 잇기 위한
실무 가이드다. 코드를 어디부터 읽고, 어떻게 실행\textbf{·}디버그\textbf{·}빌드하며,
어떤 작업에 어떤 파일을 고쳐야 하는지를 한곳에 모은다.

\subsection{코드 분석 진입점과 권장 읽기 순서}

처음 코드를 여는 동료는 ``큰 그림 $\to$ 제어 흐름 $\to$ 관심 영역''의 순서로
읽기를 권한다.

\begin{enumerate}
  \item \textbf{큰 그림}: 1장(개요)과 9장(연계 구조)으로 세 도메인과 명령 계약을 파악한다.
  \item \textbf{제어 흐름 따라가기}: 아래 두 진입점에서 호출 사슬을 손으로 따라간다.
  \item \textbf{관심 영역 심화}: 펌웨어는 4장, 웹/블록은 5\textasciitilde 7장, 3D/빌드는 8장,
        세부 시그니처는 부록 A(API Reference)를 참조한다.
\end{enumerate}

\begin{longtable}{L{2.6cm} L{12.0cm}}
\toprule
\rowcolor{tablehead}\textbf{영역} & \textbf{읽기 진입점(호출 사슬)} \\
\midrule
\endhead
펌웨어 & \code{main.py}: \code{AresRover.run()} $\to$ \code{\_read\_uart\_line()} $\to$
  \code{CommandProcessor.process()} $\to$ 개별 핸들러 $\to$ \code{robot}(\code{hardware.py}) \\
웹 & \code{main.js}: \code{main()} 초기화 $\to$ 실행 버튼 $\to$
  \code{CommandExecutor.executeWorkspace()} $\to$ \code{processBlock()} $\to$
  \code{generateCommand()} $\to$ \code{sendCommand()} $\to$ \code{bluetooth.sendData()} \\
\bottomrule
\end{longtable}

\subsection{개발 환경 구축}

\subsubsection{웹 UI}

\begin{lstlisting}[style=json, caption={웹 UI 로컬 실행}]
cd Web
python -m http.server 8000
# 브라우저에서 http://localhost:8000/main.html 열기
\end{lstlisting}

\begin{infobox}[title=주의]
개발 중에는 Web Bluetooth가 \textbf{보안 컨텍스트}(\code{http://localhost} 또는
HTTPS)에서만 동작하므로 소스를 \code{file://}로 바로 열면 BLE가 연결되지 않는다.
오프라인 단독 실행본(\code{file://})은 8장의 빌드 절차로 만들어야 한다.
\end{infobox}

\subsubsection{피코 펌웨어}

\code{Pico/}의 모든 \code{.py} 파일을 Thonny\textbf{·}\code{ampy}\textbf{·}\code{rshell}
중 하나로 라즈베리파이 피코에 업로드한다. 시리얼 디버그는 9600\,baud이며,
펌웨어는 \code{print()} 출력으로 부팅\textbf{·}모듈 초기화 상태를 보고한다.
\code{Pico/backup/}과 날짜 백업 파일은 수정 대상이 아니다.

\subsubsection{AI 모듈}

\code{AI/ai.py}는 PyTorch + Hugging Face Transformers로 로컬 LLM
(EXAONE-3.5-2.4B)을 구동하는 독립 PoC다. 별도의 파이썬 환경(GPU 권장)이
필요하며, 현재 웹 편집기와는 분리되어 있다(11장 통합 계획 참조).

\subsection{변경 영향 지도 --- 작업별 수정 파일}

자주 발생하는 변경에 대해 \emph{함께} 수정해야 하는 파일을 정리한다. 특히 웹과
펌웨어를 잇는 명령은 \textbf{양쪽}을 동시에 고쳐야 한다(9장 계약 규칙 참조).

\begin{longtable}{L{4.6cm} L{9.4cm}}
\toprule
\rowcolor{tablehead}\textbf{작업} & \textbf{수정할 파일} \\
\midrule
\endhead
새 Blockly 블록 추가 & \code{Web/blocklyconfig.js}, \code{Web/main.html}(툴박스 XML),
  \code{Web/commandexecutor.js} \\
새 펌웨어 명령 추가 & \code{Pico/process\_data.py}(\code{CommandProcessor}, 접두 순서 유의) \\
웹\textbf{·}펌웨어 공용 새 명령 & 위 두 항목을 \emph{동시에}, 응답 정책(fire-and-forget) 일치 \\
기본 시스템 설정 변경 & \code{Pico/system\_config.py}(\code{DEFAULT\_CONFIG}) \\
GPIO 핀 변경 & \code{Pico/pins.py}(또는 런타임 \code{SET\_PIN}) \\
하드웨어 모듈 추가/수정 & \code{Pico/<module>.py}, \code{Pico/hardware.py}, \code{Pico/process\_data.py} \\
BLE 타이밍/UUID 변경 & \code{Web/constants.js}, \code{Web/bluetooth.js}, \code{Pico/main.py}(UART) \\
3D 모델/시뮬레이션 추가 & \code{Web/Mesh/}, \code{Web/simulation.js}, 빌드 인라인(\code{build.*}) \\
\bottomrule
\end{longtable}

\subsection{디버깅 \textbf{·} 로깅 \textbf{·} 진단}

\begin{longtable}{L{3.0cm} L{11.0cm}}
\toprule
\rowcolor{tablehead}\textbf{수단} & \textbf{용도} \\
\midrule
\endhead
웹 통신 로그 & \code{logger.js}의 로그 UI --- 송수신 명령\textbf{·}응답을 시간순으로 확인 \\
대시보드 & \code{dashboard.html} --- 센서 텔레메트리, 수동 제어, 진단 \\
브라우저 콘솔 & 예외\textbf{·}경고, \code{state} 전역 상태 점검 \\
시리얼 모니터 & 펌웨어 \code{print()} 출력(9600\,baud) --- 부팅\textbf{·}모듈 init 실패 메시지 \\
\bottomrule
\end{longtable}

\begin{infobox}[title=흔한 실패 지점과 점검 순서]
\begin{itemize}
  \item \textbf{BLE 연결 실패}: 보안 컨텍스트(\code{localhost}/HTTPS)인가? 사용자
        제스처(버튼 클릭) 안에서 연결했는가? 장치 이름 필터\textbf{·}전원\textbf{·}거리 확인.
  \item \textbf{긴 명령이 깨짐}: 20바이트 청크 타이밍(\code{CHUNK\_DELAY})과 펌웨어
        수신 버퍼(512B)\textbf{·}타임아웃을 확인.
  \item \textbf{명령이 무시됨}: 해당 모듈이 \code{system.json}에서 비활성인가?
        \code{process\_data.py}의 접두 일치 순서가 맞는가?
  \item \textbf{응답이 멈춤}: fire-and-forget 목록이 웹\textbf{·}펌웨어 양측에서 일치하는가?
        차단형 시간 명령(\code{tFORWARD})이 루프를 점유하고 있지 않은가?
  \item \textbf{핀 변경이 안 먹힘}: \code{SET\_PIN} 후 재부팅했는가? 핀 충돌\textbf{·}예비 핀 확인.
\end{itemize}
\end{infobox}

\subsection{오프라인 빌드 \textbf{·} 배포와 기기 간 동기화}

\begin{itemize}
  \item \textbf{오프라인 빌드}: \code{build.sh}(macOS/Linux)\textbf{·}\code{build.ps1}(Windows)로
        \code{Web/}을 \code{Build/} 단일 폴더로 변환한다(8장 상세). 산출물은
        \code{file://}로 더블클릭 실행한다.
  \item \textbf{버전관리}: 저장소는 클라우드 동기화 폴더(OneDrive 등) \emph{밖}에 clone 한다.
        OneDrive ``파일 온디맨드''가 \code{.git} 객체를 손상시킨 사고가 있어, 기기 간
        동기화는 git \code{push}/\code{pull}로만 한다.
  \item \textbf{배포 사본}: 빌드 결과는 저장소에 커밋하고, 배포용 OneDrive
        \code{Build/} 폴더에도 복사해 최신으로 유지한다(\code{.git} 제외).
\end{itemize}

\subsection{공동 개발 분석 체크리스트}

새 기능을 분석\textbf{·}구현하기 전에 다음을 점검하면 회귀와 누락을 줄일 수 있다.

\begin{enumerate}
  \item 이 변경은 \emph{웹\textbf{·}펌웨어 중 어느 쪽}에 영향을 주는가? 명령 프로토콜을
        건드린다면 \textbf{양쪽}을 함께 바꿨는가?
  \item 새 명령의 \textbf{응답 정책}(fire-and-forget 여부)을 양측에서 일치시켰는가?
  \item 새 블록은 \code{blocklyconfig.js}\textbf{·}툴박스\textbf{·}\code{commandexecutor.js}
        \emph{세 곳}을 모두 갱신했는가?
  \item 새 하드웨어는 모듈 활성 플래그\textbf{·}핀 할당\textbf{·}\code{process\_data.py}
        널 검사를 갖췄는가?
  \item 긴 명령이라면 \textbf{BLE 청크\textbf{·}버퍼 한계}를 고려했는가?
  \item 실기 없이 \textbf{시뮬레이터}로도 동작이 검증되는가?
  \item 변경을 \code{Document/API\_DOCUMENTATION.md} 등 \textbf{단일 출처 문서}에 반영했는가?
\end{enumerate}

\begin{teachbox}{분석을 개발로 잇기}
\teachhead{설계 의도}
\begin{itemize}
  \item 이 가이드는 ``코드를 이해한다''와 ``코드를 바꾼다'' 사이의 간극을 줄이기 위한 것이다.
        진입점\textbf{·}영향 지도\textbf{·}체크리스트로 \emph{안전한 변경}을 돕는다.
\end{itemize}
\teachhead{핵심 질문}
\begin{itemize}
  \item 내가 맡은 작업의 ``변경 영향 지도''상 위치는 어디이며, 함께 바뀌어야 할 파일은 무엇인가?
  \item 변경을 어떻게 \emph{실행\textbf{·}시뮬\textbf{·}시리얼}로 검증할 것인가?
\end{itemize}
\teachhead{점검 요소}
\begin{itemize}
  \item 첫 작업으로, 작은 변경(예: 블록 1개 추가) 하나를 진입점부터 검증까지 끝까지
        수행하며 본 가이드의 각 단계를 실제로 밟아 보라.
\end{itemize}
\end{teachbox}
